
% https://akaedu.github.io/book/ch32.html

\chapter{正则表达式}
\BookAnchor{db.ch32}\BookAnchor{db.ch32.id2878887}%


\BookNav{\hyperlink{db.ch31s06}{上一页}\enspace\textperiodcentered\enspace \hyperlink{db.pt03}{上一级}\enspace\textperiodcentered\enspace \hyperlink{db.index}{起始页}\enspace\textperiodcentered\enspace \hyperlink{db.ch32s01}{下一页}}


% https://akaedu.github.io/book/ch32s01.html

\section{引言}
\BookAnchor{db.ch32s01}\BookAnchor{db.ch32s01.id2878891}%
以前我们用\BookCode{grep}在一个文件中找出包含某些字符串的行，比如在头文件中找出一个宏定义。其实\BookCode{grep}还可以找出\emph{符合某个模式（Pattern）\BookAnchor{db.ch32s01.id2878281}的一类字符串}。例如找出所有符合\BookCode{xxxxx@\BookSplit{}xxxx.\BookSplit{}xxx}模式的字符串（也就是email地址），要求x字符可以是字母、数字、下划线、小数点或减号，email地址的每一部分可以有一个或多个x字符，例如\BookCode{abc.\BookSplit{}d@\BookSplit{}ef.\BookSplit{}com}、\BookCode{1\_\BookSplit{}2@\BookSplit{}987-\BookSplit{}6.\BookSplit{}54}，当然符合这个模式的不全是合法的email地址，但至少可以做一次初步筛选，筛掉\BookCode{a.\BookSplit{}b}、\BookCode{c@\BookSplit{}d}等肯定不是email地址的字符串。再比如，找出所有符合\BookCode{yyy.\BookSplit{}yyy.\BookSplit{}yyy.\BookSplit{}yyy}模式的字符串（也就是IP地址），要求y是0-9的数字，IP地址的每一部分可以有1-3个y字符。



如果要用\BookCode{grep}查找一个模式，如何表示这个模式，这一类字符串，而不是一个特定的字符串呢？从这两个简单的例子可以看出，要表示一个模式至少应该包含以下信息：


\begin{itemize}

\item 字符类（Character Class）\BookAnchor{db.ch32s01.id2878960}：如上例的x和y，它们在模式中表示一个字符，但是取值范围是一类字符中的任意一个。

\item 数量限定符（Quantifier）\BookAnchor{db.ch32s01.id2878974}： 邮件地址的每一部分可以有\emph{一个或多个}x字符，IP地址的每一部分可以有\emph{1-3个}y字符

\item 各种字符类以及普通字符之间的位置关系：例如邮件地址分三部分，用普通字符\BookCode{@\BookSplit{}}和\BookCode{.\BookSplit{}}隔开，IP地址分四部分，用\BookCode{.\BookSplit{}}隔开，每一部分都可以用字符类和数量限定符描述。为了表示位置关系，还有位置限定符（Anchor）\BookAnchor{db.ch32s01.id2879017}的概念，将在下面介绍。
\end{itemize}


规定一些特殊语法表示字符类、数量限定符和位置关系，然后用这些特殊语法和普通字符一起表示一个模式，这就是正则表达式（Regular Expression）\BookAnchor{db.ch32s01.id2879032}。例如email地址的正则表达式可以写成\BookCode{[\BookSplit{}a-\BookSplit{}zA-\BookSplit{}Z0-\BookSplit{}9\_\BookSplit{}.\BookSplit{}-\BookSplit{}]\BookSplit{}+\BookSplit{}@\BookSplit{}[\BookSplit{}a-\BookSplit{}zA-\BookSplit{}Z0-\BookSplit{}9\_\BookSplit{}.\BookSplit{}-\BookSplit{}]\BookSplit{}+\BookSplit{}\textbackslash{}.\BookSplit{}[\BookSplit{}a-\BookSplit{}zA-\BookSplit{}Z0-\BookSplit{}9\_\BookSplit{}.\BookSplit{}-\BookSplit{}]\BookSplit{}+\BookSplit{}}，IP地址的正则表达式可以写成\BookCode{[\BookSplit{}0-\BookSplit{}9]\BookSplit{}\{1,\BookSplit{}3\}\textbackslash{}.\BookSplit{}[\BookSplit{}0-\BookSplit{}9]\BookSplit{}\{1,\BookSplit{}3\}\textbackslash{}.\BookSplit{}[\BookSplit{}0-\BookSplit{}9]\BookSplit{}\{1,\BookSplit{}3\}\textbackslash{}.\BookSplit{}[\BookSplit{}0-\BookSplit{}9]\BookSplit{}\{1,\BookSplit{}3\}}。下一节介绍正则表达式的语法，我们先看看正则表达式在\BookCode{grep}中怎么用。例如有这样一个文本文件\BookCode{testfile}：


\BookPlain{code/ch32s01/01.txt}{6}{4}


查找其中包含IP地址的行：


\BookTerminal{code/ch32s01/02.txt}{5}{3}


\BookCode{egrep}相当于\BookCode{grep -\BookSplit{}E}，表示采用Extended正则表达式语法。\BookCode{grep}的正则表达式有Basic和Extended两种规范，它们之间的区别下一节再解释。另外还有\BookCode{fgrep}命令，相当于\BookCode{grep -\BookSplit{}F}，表示只搜索固定字符串而不搜索正则表达式模式，不会按正则表达式的语法解释后面的参数。



注意正则表达式参数用单引号括起来了，因为正则表达式中用到的很多特殊字符在Shell中也有特殊含义（例如\textbackslash{}），只有用单引号括起来才能保证这些字符原封不动地传给\BookCode{grep}命令，而不会被Shell解释掉。



\BookCode{192.\BookSplit{}168.\BookSplit{}1.\BookSplit{}1}符合上述模式，由三个\BookCode{.\BookSplit{}}隔开的四段组成，每段都是1到3个数字，所以这一行被找出来了，可为什么\BookCode{1234.\BookSplit{}234.\BookSplit{}04.\BookSplit{}5678}也被找出来了呢？因为\BookCode{grep}找的是\emph{包含}某一模式的行，这一行包含一个符合模式的字符串\BookCode{234.\BookSplit{}234.\BookSplit{}04.\BookSplit{}567}。相反，\BookCode{123.\BookSplit{}4234.\BookSplit{}045.\BookSplit{}678}这一行不包含符合模式的字符串，所以不会被找出来。



\BookCode{grep}是一种查找过滤工具，正则表达式在\BookCode{grep}中用来查找符合模式的字符串。其实正则表达式还有一个重要的应用是验证用户输入是否合法，例如用户通过网页表单提交自己的email地址，就需要用程序验证一下是不是合法的email地址，这个工作可以在网页的Javascript中做，也可以在网站后台的程序中做，例如PHP、Perl、Python、Ruby、Java或C，所有这些语言都支持正则表达式，可以说，目前不支持正则表达式的编程语言实在很少见。除了编程语言之外，很多UNIX命令和工具也都支持正则表达式，例如grep、vi、sed、awk、emacs等等。“正则表达式”就像“变量”一样，它是一个广泛的概念，而不是某一种工具或编程语言的特性。

\BookNav{\hyperlink{db.ch32}{上一页}\enspace\textperiodcentered\enspace \hyperlink{db.ch32}{上一级}\enspace\textperiodcentered\enspace \hyperlink{db.index}{起始页}\enspace\textperiodcentered\enspace \hyperlink{db.ch32s02}{下一页}}


% https://akaedu.github.io/book/ch32s02.html

\section{基本语法}
\BookAnchor{db.ch32s02}\BookAnchor{db.ch32s02.id2879223}%
我们知道C的变量和Shell脚本变量的定义和使用方法很不相同，表达能力也不相同，C的变量有各种类型，而Shell脚本变量都是字符串。同样道理，各种工具和编程语言所使用的正则表达式规范的语法并不相同，表达能力也各不相同，有的正则表达式规范引入很多扩展，能表达更复杂的模式，但各种正则表达式规范的基本概念都是相通的。本节介绍\BookCode{egrep(\BookSplit{}1)\BookSplit{}}所使用的正则表达式，它大致上符合POSIX正则表达式规范，详见\BookCode{regex(\BookSplit{}7)\BookSplit{}}（看这个man page对你的英文绝对是很好的锻炼）。希望读者仿照上一节的例子，一边学习语法，一边用\BookCode{egrep}命令做实验。


\begingroup
\setlength{\tabcolsep}{6pt}
\small
\sloppy
\renewcommand{\arraystretch}{1.16}
\setlength{\parskip}{0pt}
\begin{longtable}{L{\dimexpr 0.156447\linewidth-2\tabcolsep\relax}L{\dimexpr 0.398608\linewidth-2\tabcolsep\relax}L{\dimexpr 0.444945\linewidth-2\tabcolsep\relax}}
\BookTableTitle{字符类}\BookTableLabel{\BookAnchor{db.ch32s02.id2879268}}
\toprule
{\bfseries 字符} & {\bfseries 含义} & {\bfseries 举例} \\
\midrule
\endfirsthead
\toprule
{\bfseries 字符} & {\bfseries 含义} & {\bfseries 举例} \\
\midrule
\endhead
\BookCode{.\BookSplit{}} & 匹配任意一个字符 & \BookCode{abc.\BookSplit{}}可以匹配\BookCode{abcd}、\BookCode{abc9}等 \\
\addlinespace[2pt]
\BookCode{[\BookSplit{}]\BookSplit{}} & 匹配括号中的任意一个字符 & \BookCode{[\BookSplit{}abc]\BookSplit{}d}可以匹配\BookCode{ad}、\BookCode{bd}或\BookCode{cd} \\
\addlinespace[2pt]
\BookCode{-\BookSplit{}} & 在\BookCode{[\BookSplit{}]\BookSplit{}}括号内表示字符范围 & \BookCode{[\BookSplit{}0-\BookSplit{}9a-\BookSplit{}fA-\BookSplit{}F]\BookSplit{}}\BookCode{}可以匹配一位十六进制数字 \\
\addlinespace[2pt]
\BookCode{\textasciicircum{}} & 位于\BookCode{[\BookSplit{}]\BookSplit{}}括号内的开头，匹配除括号中的字符之外的任意一个字符 & \BookCode{[\BookSplit{}\textasciicircum{}xy]\BookSplit{}}匹配除\BookCode{xy}之外的任一字符，因此\BookCode{[\BookSplit{}\textasciicircum{}xy]\BookSplit{}1}可以匹配\BookCode{a1}、\BookCode{b1}但不匹配\BookCode{x1}、\BookCode{y1} \\
\addlinespace[2pt]
\BookCode{[\BookSplit{}[\BookSplit{}:\BookSplit{}xxx:\BookSplit{}]\BookSplit{}]\BookSplit{}} & \BookCode{grep}工具预定义的一些命名字符类 & \BookCode{[\BookSplit{}[\BookSplit{}:\BookSplit{}alpha:\BookSplit{}]\BookSplit{}]\BookSplit{}}匹配一个字母，\BookCode{[\BookSplit{}[\BookSplit{}:\BookSplit{}digit:\BookSplit{}]\BookSplit{}]\BookSplit{}}匹配一个数字 \\
\bottomrule
\end{longtable}
\endgroup

\begingroup
\setlength{\tabcolsep}{6pt}
\small
\sloppy
\renewcommand{\arraystretch}{1.16}
\setlength{\parskip}{0pt}
\begin{longtable}{L{\dimexpr 0.103853\linewidth-2\tabcolsep\relax}L{\dimexpr 0.347508\linewidth-2\tabcolsep\relax}L{\dimexpr 0.548639\linewidth-2\tabcolsep\relax}}
\BookTableTitle{数量限定符}\BookTableLabel{\BookAnchor{db.ch32s02.id2879472}}
\toprule
{\bfseries 字符} & {\bfseries 含义} & {\bfseries 举例} \\
\midrule
\endfirsthead
\toprule
{\bfseries 字符} & {\bfseries 含义} & {\bfseries 举例} \\
\midrule
\endhead
\BookCode{?} & 紧跟在它前面的单元应匹配零次或一次 & \BookCode{[\BookSplit{}0-\BookSplit{}9]\BookSplit{}?\textbackslash{}.\BookSplit{}[\BookSplit{}0-\BookSplit{}9]\BookSplit{}}匹配\BookCode{0.\BookSplit{}0}、\BookCode{2.\BookSplit{}3}、\BookCode{.\BookSplit{}5}等，由于\BookCode{.\BookSplit{}}在正则表达式中是一个特殊字符，所以需要用\BookCode{\textbackslash{}}转义一下，取字面值 \\
\addlinespace[2pt]
\BookCode{+\BookSplit{}} & 紧跟在它前面的单元应匹配一次或多次 & \BookCode{[\BookSplit{}a-\BookSplit{}zA-\BookSplit{}Z0-\BookSplit{}9\_\BookSplit{}.\BookSplit{}-\BookSplit{}]\BookSplit{}+\BookSplit{}@\BookSplit{}[\BookSplit{}a-\BookSplit{}zA-\BookSplit{}Z0-\BookSplit{}9\_\BookSplit{}.\BookSplit{}-\BookSplit{}]\BookSplit{}+\BookSplit{}\textbackslash{}.\BookSplit{}[\BookSplit{}a-\BookSplit{}zA-\BookSplit{}Z0-\BookSplit{}9\_\BookSplit{}.\BookSplit{}-\BookSplit{}]\BookSplit{}+\BookSplit{}}匹配email地址 \\
\addlinespace[2pt]
\BookCode{*} & 紧跟在它前面的单元应匹配零次或多次 & \BookCode{[\BookSplit{}0-\BookSplit{}9]\BookSplit{}[\BookSplit{}0-\BookSplit{}9]\BookSplit{}*}匹配至少一位数字，等价于\BookCode{[\BookSplit{}0-\BookSplit{}9]\BookSplit{}+\BookSplit{}}，\BookCode{[\BookSplit{}a-\BookSplit{}zA-\BookSplit{}Z\_\BookSplit{}]\BookSplit{}+\BookSplit{}[\BookSplit{}a-\BookSplit{}zA-\BookSplit{}Z\_\BookSplit{}0-\BookSplit{}9]\BookSplit{}*}匹配C语言的标识符 \\
\addlinespace[2pt]
\BookCode{\{N\}} & 紧跟在它前面的单元应精确匹配\BookCode{N次} & \BookCode{[\BookSplit{}1-\BookSplit{}9]\BookSplit{}[\BookSplit{}0-\BookSplit{}9]\BookSplit{}\{2\}}匹配从\BookCode{100}到\BookCode{999}的整数 \\
\addlinespace[2pt]
\BookCode{\{N,\BookSplit{}\}} & 紧跟在它前面的单元应匹配至少\BookCode{N}次 & \BookCode{[\BookSplit{}1-\BookSplit{}9]\BookSplit{}[\BookSplit{}0-\BookSplit{}9]\BookSplit{}\{2,\BookSplit{}\}}匹配三位以上（含三位）的整数 \\
\addlinespace[2pt]
\BookCode{\{,\BookSplit{}M\}} & 紧跟在它前面的单元应匹配最多\BookCode{M}次 & \BookCode{[\BookSplit{}0-\BookSplit{}9]\BookSplit{}\{,\BookSplit{}1\}}相当于\BookCode{[\BookSplit{}0-\BookSplit{}9]\BookSplit{}?} \\
\addlinespace[2pt]
\BookCode{\{N,\BookSplit{}M\}} & 紧跟在它前面的单元应匹配至少\BookCode{N}次，最多\BookCode{M}次 & \BookCode{[\BookSplit{}0-\BookSplit{}9]\BookSplit{}\{1,\BookSplit{}3\}\textbackslash{}.\BookSplit{}[\BookSplit{}0-\BookSplit{}9]\BookSplit{}\{1,\BookSplit{}3\}\textbackslash{}.\BookSplit{}[\BookSplit{}0-\BookSplit{}9]\BookSplit{}\{1,\BookSplit{}3\}\textbackslash{}.\BookSplit{}[\BookSplit{}0-\BookSplit{}9]\BookSplit{}\{1,\BookSplit{}3\}}匹配IP地址 \\
\bottomrule
\end{longtable}
\endgroup


再次注意\BookCode{grep}找的是包含某一模式的行，而不是完全匹配某一模式的行。再举个例子，如果文本文件的内容是


\BookPlain{code/ch32s02/01.txt}{5}{3}


查找\BookCode{a*}这个模式的结果是三行都被找出来了


\BookTerminal{code/ch32s02/02.txt}{6}{4}


\BookCode{a*}匹配0个或多个\BookCode{a}，而第三行包含0个\BookCode{a}，所以也包含了这一模式。单独用\BookCode{a*}这样的正则表达式做查找没什么意义，一般是把\BookCode{a*}作为正则表达式的一部分来用。


\begingroup
\setlength{\tabcolsep}{6pt}
\small
\sloppy
\renewcommand{\arraystretch}{1.16}
\setlength{\parskip}{0pt}
\begin{longtable}{L{\dimexpr 0.083845\linewidth-2\tabcolsep\relax}L{\dimexpr 0.335380\linewidth-2\tabcolsep\relax}L{\dimexpr 0.580775\linewidth-2\tabcolsep\relax}}
\BookTableTitle{位置限定符}\BookTableLabel{\BookAnchor{db.ch32s02.id2879776}}
\toprule
{\bfseries 字符} & {\bfseries 含义} & {\bfseries 举例} \\
\midrule
\endfirsthead
\toprule
{\bfseries 字符} & {\bfseries 含义} & {\bfseries 举例} \\
\midrule
\endhead
\BookCode{\textasciicircum{}} & 匹配行首的位置 & \BookCode{\textasciicircum{}Content}匹配位于一行开头的\BookCode{Content} \\
\addlinespace[2pt]
\BookCode{\$} & 匹配行末的位置 & \BookCode{;\BookSplit{}\$}匹配位于一行结尾的\BookCode{;\BookSplit{}}号，\BookCode{\textasciicircum{}\$}匹配空行 \\
\addlinespace[2pt]
\BookCode{\textbackslash{}\textless{}} & 匹配单词开头的位置 & \BookCode{\textbackslash{}\textless{}th}匹配\BookCode{.\BookSplit{}.\BookSplit{}.\BookSplit{} this}，但不匹配\BookCode{ethernet}、\BookCode{tenth} \\
\addlinespace[2pt]
\BookCode{\textbackslash{}\textgreater{}} & 匹配单词结尾的位置 & \BookCode{p\textbackslash{}\textgreater{}}匹配\BookCode{leap .\BookSplit{}.\BookSplit{}.\BookSplit{}}，但不匹配\BookCode{parent}、\BookCode{sleepy} \\
\addlinespace[2pt]
\BookCode{\textbackslash{}b} & 匹配单词开头或结尾的位置 & \BookCode{\textbackslash{}bat\textbackslash{}b}匹配\BookCode{.\BookSplit{}.\BookSplit{}.\BookSplit{} at .\BookSplit{}.\BookSplit{}.\BookSplit{}}，但不匹配\BookCode{cat}、\BookCode{atexit}、\BookCode{batch} \\
\addlinespace[2pt]
\BookCode{\textbackslash{}B} & 匹配非单词开头和结尾的位置 & \BookCode{\textbackslash{}Bat\textbackslash{}B}匹配\BookCode{battery}，但不匹配\BookCode{.\BookSplit{}.\BookSplit{}.\BookSplit{} attend}、\BookCode{hat .\BookSplit{}.\BookSplit{}.\BookSplit{}} \\
\bottomrule
\end{longtable}
\endgroup


位置限定符可以帮助\BookCode{grep}更准确地查找，例如上一节我们用\BookCode{[\BookSplit{}0-\BookSplit{}9]\BookSplit{}\{1,\BookSplit{}3\}\textbackslash{}.\BookSplit{}[\BookSplit{}0-\BookSplit{}9]\BookSplit{}\{1,\BookSplit{}3\}\textbackslash{}.\BookSplit{}[\BookSplit{}0-\BookSplit{}9]\BookSplit{}\{1,\BookSplit{}3\}\textbackslash{}.\BookSplit{}[\BookSplit{}0-\BookSplit{}9]\BookSplit{}\{1,\BookSplit{}3\}}查找IP地址，找到这两行


\BookTerminal{code/ch32s02/03.txt}{4}{2}


如果用\BookCode{\textasciicircum{}[\BookSplit{}0-\BookSplit{}9]\BookSplit{}\{1,\BookSplit{}3\}\textbackslash{}.\BookSplit{}[\BookSplit{}0-\BookSplit{}9]\BookSplit{}\{1,\BookSplit{}3\}\textbackslash{}.\BookSplit{}[\BookSplit{}0-\BookSplit{}9]\BookSplit{}\{1,\BookSplit{}3\}\textbackslash{}.\BookSplit{}[\BookSplit{}0-\BookSplit{}9]\BookSplit{}\{1,\BookSplit{}3\}\$}查找，就可以把\BookCode{1234.\BookSplit{}234.\BookSplit{}04.\BookSplit{}5678}这一行过滤掉了。


\begingroup
\setlength{\tabcolsep}{6pt}
\small
\sloppy
\renewcommand{\arraystretch}{1.16}
\setlength{\parskip}{0pt}
\begin{longtable}{L{\dimexpr 0.083845\linewidth-2\tabcolsep\relax}L{\dimexpr 0.436167\linewidth-2\tabcolsep\relax}L{\dimexpr 0.479988\linewidth-2\tabcolsep\relax}}
\BookTableTitle{其它特殊字符}\BookTableLabel{\BookAnchor{db.ch32s02.id2880045}}
\toprule
{\bfseries 字符} & {\bfseries 含义} & {\bfseries 举例} \\
\midrule
\endfirsthead
\toprule
{\bfseries 字符} & {\bfseries 含义} & {\bfseries 举例} \\
\midrule
\endhead
\BookCode{\textbackslash{}} & 转义字符，普通字符转义为特殊字符，特殊字符转义为普通字符 & 普通字符\BookCode{\textless{}}写成\BookCode{\textbackslash{}\textless{}}表示单词开头的位置，特殊字符\BookCode{.\BookSplit{}}写成\BookCode{\textbackslash{}.\BookSplit{}}以及\BookCode{\textbackslash{}}写成\BookCode{\textbackslash{}\textbackslash{}}就当作普通字符来匹配 \\
\addlinespace[2pt]
\BookCode{(\BookSplit{})\BookSplit{}} & 将正则表达式的一部分括起来组成一个单元，可以对整个单元使用数量限定符 & \BookCode{(\BookSplit{}[\BookSplit{}0-\BookSplit{}9]\BookSplit{}\{1,\BookSplit{}3\}\textbackslash{}.\BookSplit{})\BookSplit{}\{3\}[\BookSplit{}0-\BookSplit{}9]\BookSplit{}\{1,\BookSplit{}3\}}匹配IP地址 \\
\addlinespace[2pt]
\BookCode{|\BookSplit{}} & 连接两个子表达式，表示或的关系 & \BookCode{n(\BookSplit{}o|\BookSplit{}either)\BookSplit{}}匹配\BookCode{no}或\BookCode{neither} \\
\bottomrule
\end{longtable}
\endgroup


以上介绍的是\BookCode{grep}正则表达式的Extended规范，Basic规范也有这些语法，只是字符\BookCode{?+\BookSplit{}\{\}|\BookSplit{}(\BookSplit{})\BookSplit{}}应解释为普通字符，要表示上述特殊含义则需要加\BookCode{\textbackslash{}}转义。如果用\BookCode{grep}而不是\BookCode{egrep}，并且不加\BookCode{-\BookSplit{}E}参数，则应该遵照Basic规范来写正则表达式。

\BookNav{\hyperlink{db.ch32s01}{上一页}\enspace\textperiodcentered\enspace \hyperlink{db.ch32}{上一级}\enspace\textperiodcentered\enspace \hyperlink{db.index}{起始页}\enspace\textperiodcentered\enspace \hyperlink{db.ch32s03}{下一页}}


% https://akaedu.github.io/book/ch32s03.html

\section{sed}
\BookAnchor{db.ch32s03}\BookAnchor{db.ch32s03.id2880212}%
\BookCode{sed}意为流编辑器（Stream Editor）\BookAnchor{db.ch32s03.id2880224}，在Shell脚本和Makefile中作为过滤器使用非常普遍，也就是把前一个程序的输出引入sed的输入，经过一系列编辑命令转换为另一种格式输出。\BookCode{sed}和\BookCode{vi}都源于早期UNIX的\BookCode{ed}工具，所以很多\BookCode{sed}命令和\BookCode{vi}的末行命令是相同的。



\BookCode{sed}命令行的基本格式为


\BookPlain{code/ch32s03/01.txt}{4}{2}


\BookCode{sed}处理的文件既可以由标准输入重定向得到，也可以当命令行参数传入，命令行参数可以一次传入多个文件，\BookCode{sed}会依次处理。\BookCode{sed}的编辑命令可以直接当命令行参数传入，也可以写成一个脚本文件然后用\BookCode{-\BookSplit{}f}参数指定，编辑命令的格式为


\BookPlain{code/ch32s03/02.txt}{3}{1}


其中\BookCode{pattern}是正则表达式，\BookCode{action}是编辑操作。\BookCode{sed}程序一行一行读出待处理文件，如果某一行与\BookCode{pattern}匹配，则执行相应的\BookCode{action}，如果一条命令没有\BookCode{pattern}而只有\BookCode{action}，这个\BookCode{action}将作用于待处理文件的每一行。


\begingroup
\setlength{\tabcolsep}{6pt}
\small
\sloppy
\renewcommand{\arraystretch}{1.16}
\setlength{\parskip}{0pt}
\begin{longtable}{L{\dimexpr 0.411693\linewidth-2\tabcolsep\relax}L{\dimexpr 0.588307\linewidth-2\tabcolsep\relax}}
\BookTableTitle{常用的sed命令}\BookTableLabel{\BookAnchor{db.ch32s03.id2880365}}
\toprule
\BookCode{/\BookSplit{}pattern/\BookSplit{}p} & 打印匹配\BookCode{pattern}的行 \\
\addlinespace[2pt]
\BookCode{/\BookSplit{}pattern/\BookSplit{}d} & 删除匹配\BookCode{pattern}的行 \\
\addlinespace[2pt]
\BookCode{/\BookSplit{}pattern/\BookSplit{}s/\BookSplit{}pattern1/\BookSplit{}pattern2/\BookSplit{}} & 查找符合\BookCode{pattern}的行，将该行第一个匹配\BookCode{pattern1}的字符串替换为\BookCode{pattern2} \\
\addlinespace[2pt]
\BookCode{/\BookSplit{}pattern/\BookSplit{}s/\BookSplit{}pattern1/\BookSplit{}pattern2/\BookSplit{}g} & 查找符合\BookCode{pattern}的行，将该行所有匹配\BookCode{pattern1}的字符串替换为\BookCode{pattern2} \\
\bottomrule
\end{longtable}
\endgroup


使用\BookCode{p}命令需要注意，\BookCode{sed}是把待处理文件的内容连同处理结果一起输出到标准输出的，因此\BookCode{p}命令表示除了把文件内容打印出来之外还额外打印一遍匹配\BookCode{pattern}的行。比如一个文件\BookCode{testfile}的内容是


\BookPlain{code/ch32s03/03.txt}{5}{3}


打印其中包含\BookCode{abc}的行


\BookTerminal{code/ch32s03/04.txt}{7}{5}


要想只输出处理结果，应加上\BookCode{-\BookSplit{}n}选项，这种用法相当于\BookCode{grep}命令


\BookTerminal{code/ch32s03/05.txt}{4}{2}


使用\BookCode{d}命令就不需要\BookCode{-\BookSplit{}n}参数了，比如删除含有\BookCode{abc}的行


\BookTerminal{code/ch32s03/06.txt}{5}{3}


注意，\BookCode{sed}命令不会修改原文件，删除命令只表示某些行不打印输出，而不是从原文件中删去。



使用查找替换命令时，可以把匹配\BookCode{pattern1}的字符串复制到\BookCode{pattern2}中，比如：


\BookTerminal{code/ch32s03/07.txt}{6}{4}


\BookCode{pattern2}中的\BookCode{\&}表示原文件的当前行中与\BookCode{pattern1}相匹配的字符串，再比如：


\BookTerminal{code/ch32s03/08.txt}{6}{4}


\BookCode{pattern2}中的\BookCode{\textbackslash{}1}表示与\BookCode{pattern1}的第一个\BookCode{(\BookSplit{})\BookSplit{}}括号相匹配的内容，\BookCode{\textbackslash{}2}表示与\BookCode{pattern1}的第二个\BookCode{(\BookSplit{})\BookSplit{}}括号相匹配的内容。\BookCode{sed}默认使用Basic正则表达式规范，如果指定了\BookCode{-\BookSplit{}r}选项则使用Extended规范，那么\BookCode{(\BookSplit{})\BookSplit{}}括号就不必转义了。



如果\BookCode{testfile}的内容是


\BookPlain{code/ch32s03/09.txt}{4}{2}


现在要去掉所有的HTML标签，使输出结果为


\BookTerminal{code/ch32s03/10.txt}{4}{2}


怎么做呢？如果用下面的命令


\BookTerminal{code/ch32s03/11.txt}{6}{4}


结果是两个空行，把所有字符都过滤掉了。这是因为，正则表达式中的数量限定符会匹配尽可能长的字符串，这称为贪心的（Greedy）\BookAnchor{db.ch32s03.id2880724}\BookAnchor{db.ch32s03.id2880730}\footnote[39]{\BookAnchor{db.ch32s03.ftn.id2880730}\BookFootnoteBack{db.ch32s03.id2880730}有些正则表达式规范支持Non-greedy的数量限定符，匹配尽可能短的字符串，例如在Python中\BookCode{*?}和\BookCode{*}一样表示0个或任意多个，但前者是Non-greedy的。}。比如\BookCode{sed}在处理第一行时，\BookCode{\textless{}.\BookSplit{}*\textgreater{}}匹配的并不是\BookCode{\textless{}html\textgreater{}}或\BookCode{\textless{}head\textgreater{}}这样的标签，而是


\BookPlain{code/ch32s03/12.txt}{3}{1}


这样一整行，因为这一行开头是\BookCode{\textless{}}，中间是若干个任意字符，末尾是\BookCode{\textgreater{}}。那么这条命令怎么改才对呢？留给读者思考。

\BookNav{\hyperlink{db.ch32s02}{上一页}\enspace\textperiodcentered\enspace \hyperlink{db.ch32}{上一级}\enspace\textperiodcentered\enspace \hyperlink{db.index}{起始页}\enspace\textperiodcentered\enspace \hyperlink{db.ch32s04}{下一页}}


% https://akaedu.github.io/book/ch32s04.html

\section{awk}
\BookAnchor{db.ch32s04}\BookAnchor{db.ch32s04.id2880797}%
\BookCode{sed}以行为单位处理文件，\BookCode{awk}比\BookCode{sed}强的地方在于不仅能以行为单位还能以列为单位处理文件。\BookCode{awk}缺省的行分隔符是换行，缺省的列分隔符是连续的空格和Tab，但是行分隔符和列分隔符都可以自定义，比如\BookCode{/\BookSplit{}etc/\BookSplit{}passwd}文件的每一行有若干个字段，字段之间以\BookCode{:\BookSplit{}}分隔，就可以重新定义\BookCode{awk}的列分隔符为\BookCode{:\BookSplit{}}并以列为单位处理这个文件。\BookCode{awk}实际上是一门很复杂的脚本语言，还有像C语言一样的分支和循环结构，但是基本用法和\BookCode{sed}类似，\BookCode{awk}命令行的基本形式为：


\BookPlain{code/ch32s04/01.txt}{4}{2}


和\BookCode{sed}一样，\BookCode{awk}处理的文件既可以由标准输入重定向得到，也可以当命令行参数传入，编辑命令可以直接当命令行参数传入，也可以用\BookCode{-\BookSplit{}f}参数指定一个脚本文件，编辑命令的格式为：


\BookPlain{code/ch32s04/02.txt}{4}{2}


和\BookCode{sed}类似，\BookCode{pattern}是正则表达式，\BookCode{actions}是一系列操作。\BookCode{awk}程序一行一行读出待处理文件，如果某一行与\BookCode{pattern}匹配，或者满足\BookCode{condition}条件，则执行相应的\BookCode{actions}，如果一条\BookCode{awk}命令只有\BookCode{actions}部分，则\BookCode{actions}作用于待处理文件的每一行。比如文件\BookCode{testfile}的内容表示某商店的库存量：


\BookPlain{code/ch32s04/03.txt}{5}{3}


打印每一行的第二列:


\BookTerminal{code/ch32s04/04.txt}{6}{4}


自动变量\BookCode{\$1}、\BookCode{\$2}分别表示第一列、第二列等，类似于Shell脚本的位置参数，而\BookCode{\$0}表示整个当前行。再比如，如果某种产品的库存量低于75则在行末标注需要订货：


\BookTerminal{code/ch32s04/05.txt}{6}{4}


可见\BookCode{awk}也有和C语言非常相似的\BookCode{printf}函数。\BookCode{awk}命令的\BookCode{condition}部分还可以是两个特殊的\BookCode{condition}－\BookCode{BEGIN}和\BookCode{END}，对于每个待处理文件，\BookCode{BEGIN}后面的\BookCode{actions}在处理整个文件之前执行一次，\BookCode{END}后面的\BookCode{actions}在整个文件处理完之后执行一次。



\BookCode{awk}命令可以像C语言一样使用变量（但不需要定义变量），比如统计一个文件中的空行数


\BookTerminal{code/ch32s04/06.txt}{3}{1}


就像Shell的环境变量一样，有些\BookCode{awk}变量是预定义的有特殊含义的：


\begingroup
\setlength{\tabcolsep}{6pt}
\small
\sloppy
\renewcommand{\arraystretch}{1.16}
\setlength{\parskip}{0pt}
\begin{longtable}{L{\dimexpr 0.190889\linewidth-2\tabcolsep\relax}L{\dimexpr 0.809111\linewidth-2\tabcolsep\relax}}
\BookTableTitle{awk常用的内建变量}\BookTableLabel{\BookAnchor{db.ch32s04.id2881116}}
\toprule
FILENAME & 当前输入文件的文件名，该变量是只读的 \\
\addlinespace[2pt]
NR & 当前行的行号，该变量是只读的，\BookCode{R}代表record \\
\addlinespace[2pt]
NF & 当前行所拥有的列数，该变量是只读的，\BookCode{F}代表field \\
\addlinespace[2pt]
OFS & 输出格式的列分隔符，缺省是空格 \\
\addlinespace[2pt]
FS & 输入文件的列分融符，缺省是连续的空格和Tab \\
\addlinespace[2pt]
ORS & 输出格式的行分隔符，缺省是换行符 \\
\addlinespace[2pt]
RS & 输入文件的行分隔符，缺省是换行符 \\
\bottomrule
\end{longtable}
\endgroup


例如打印系统中的用户帐号列表


\BookTerminal{code/ch32s04/07.txt}{3}{1}


\BookCode{awk}还可以像C语言一样使用\BookCode{if}/\BookCode{else}、\BookCode{while}、\BookCode{for}控制结构，此处从略。

\BookNav{\hyperlink{db.ch32s03}{上一页}\enspace\textperiodcentered\enspace \hyperlink{db.ch32}{上一级}\enspace\textperiodcentered\enspace \hyperlink{db.index}{起始页}\enspace\textperiodcentered\enspace \hyperlink{db.ch32s05}{下一页}}


% https://akaedu.github.io/book/ch32s05.html

\section{练习：在C语言中使用正则表达式}
\BookAnchor{db.ch32s05}\BookAnchor{db.ch32s05.id2881236}%
POSIX规定了正则表达式的C语言库函数，详见\BookCode{regex(\BookSplit{}3)\BookSplit{}}。我们已经学习了很多C语言库函数的用法，读者应该具备自己看懂man手册的能力了。本章介绍了正则表达式在\BookCode{grep}、\BookCode{sed}、\BookCode{awk}中的用法，学习要能够举一反三，请读者根据\BookCode{regex(\BookSplit{}3)\BookSplit{}}自己总结正则表达式在C语言中的用法，写一些简单的程序，例如验证用户输入的IP地址或email地址格式是否正确。

\BookNav{\hyperlink{db.ch32s04}{上一页}\enspace\textperiodcentered\enspace \hyperlink{db.ch32}{上一级}\enspace\textperiodcentered\enspace \hyperlink{db.index}{起始页}\enspace\textperiodcentered\enspace \hyperlink{db.ch33}{下一页}}
